% Part: first-order-logic
% Chapter: syntax-and-semantics
% Section: terms-formulas

\documentclass[../../../include/open-logic-section]{subfiles}

\begin{document}

\olfileid{fol}{syn}{frm}

\olsection{الحدود وال\printtoken{P}{formula}}

إذا تعيّنت لغة من الرتبة الأولى~$\Lang L$، تعيّن المعجم الذي تُبنى
منه تعبيرات~$\Lang L$. ونخصّ من هذه التعبيرات هنا
\emph{الحدود} و\emph{!!{formula}s} بالتعريف.

\begin{defn}[الحدود]
\ollabel{defn:terms}
أما مجموعة \emph{الحدود}~$\Trm[L]$ في~$\Lang L$، فحدّها الاستقرائي
هذه الشروط:
\begin{enumerate}
\item كل !!{variable} حد.
\item كل !!{constant} في~$\Lang L$ حد.
\item إذا كانت $f$ !!{function} ذات $n$ من المحال، وكانت
  $t_1$، و\dots، و$t_n$ حدودًا، فإن
  $\Atom{f}{t_1, \ldots, t_n}$ حد.
\tagitem{limitClause}{
  لا شيء غير ذلك حد.}{}
\end{enumerate}
والحد الذي لا يحتوي على !!{variable}s \emph{حد مغلق}.
\end{defn}

\begin{explain}
قد أفردنا ال!!{constant}s بصنف من الرموز حين حدّدنا اللغة والحدود.
ولنا وجه آخر، وهو عدّها !!{function}s صفرية المحل؛ وعندئذ يستغنى
عن البند الثاني من تعريف الحدود. وليس علينا إلا أن نحمل $\Atom{f}{t_1, \ldots, t_n}$ على أنه $f$ وحده إذا
كان $n = 0$.
\end{explain}

\begin{defn}[الصيغ]
\ollabel{defn:formulas}
وأما مجموعة \emph{!!{formula}s}~$\Frm[L]$ في اللغة~$\Lang L$،
فتتعيّن استقرائيًا بالشروط الآتية:
\begin{enumerate}
\tagitem{prvFalse}{$\lfalse$ هي !!{formula} ذرية.}{}

\tagitem{prvTrue}{$\ltrue$ هي !!{formula} ذرية.}{}

\item إذا كان $R$ !!{predicate} ذا $n$ من المحال في~$\Lang L$، وكانت
  $t_1$، و\dots، و$t_n$ حدودًا في~$\Lang L$، فإن
  $\Atom{R}{t_1,\ldots, t_n}$ !!{formula} ذرية.

\item إذا كان $t_1$ و$t_2$ حدين في~$\Lang L$، فإن
  $\Atom{\eq}{t_1, t_2}$ !!{formula} ذرية.

\tagitem{prvNot}{إذا كانت $!A$ !!a{formula}، فإن $\lnot !A$
  !!a{formula}.}{}

\tagitem{prvAnd}{إذا كانت $!A$ و$!B$ !!{formula}s، فإن $(!A \land
  !B)$ !!a{formula}.}{}

\tagitem{prvOr}{إذا كانت $!A$ و$!B$ !!{formula}s، فإن $(!A \lor !B)$
  !!a{formula}.}{}

\tagitem{prvIf}{إذا كانت $!A$ و$!B$ !!{formula}s، فإن $(!A \lif !B)$
  !!a{formula}.}{}

\tagitem{prvIff}{إذا كانت $!A$ و$!B$ !!{formula}s، فإن
  $(!A \liff !B)$ !!a{formula}.}{}

\tagitem{prvAll}{إذا كانت $!A$ !!a{formula} وكان $x$ !!a{variable}،
  فإن $\lforall[x][!A]$ !!a{formula}.}{}

\tagitem{prvEx}{إذا كانت $!A$ !!a{formula} وكان $x$ !!a{variable}،
  فإن $\lexists[x][!A]$ !!a{formula}.}{}

\tagitem{limitClause}{لا شيء غير ذلك !!a{formula}.}{}
\end{enumerate}
\end{defn}

\begin{explain}
تعريف الحدود وتعريف ال!!{formula}s كلاهما من
\emph{التعريفات الاستقرائية}. ومعنى ذلك أن بناء مجموعة
ال!!{formula}s يجري على مراحل لا تتناهى. نبدأ بجعل الصيغ الذرية
كلها صيغًا، وهو ما تقرره الحالات الأولى من التعريف، أعني حالات
\iftag{prvTrue}{$\ltrue$، }{}%
\iftag{prvFalse}{$\lfalse$، }{}%
$\Atom{R}{t_1,\dots,t_n}$ و$\Atom{\eq}{t_1,t_2}$. فكل
!!{formula} بهذه الصورة هي المرادة باسم «!!{formula} ذرية».

وأما سائر حالات التعريف، فقواعد يُنشأ بها الجديد من ال!!{formula}s
من ال!!{formula}s التي سبق إنشاؤها. ففي المرحلة الثانية نطبّقها
فننشئ !!{formula}s من !!{formula}s ذرية؛ ثم نجمع في الثالثة ما
كان ذريًا وما نشأ في الثانية، ونبني منه صيغًا أخرى، ونجري بعد ذلك
على هذا النهج. فما حصل في مرحلة من هذه المراحل فهو !!{formula}،
وما لم يحصل في شيء منها فليس كذلك.
\end{explain}

ومن اصطلاح الكتابة أن نجعل $\eq$ بين حجتيها، مع إسقاط القوسين:
فـ~$\eq[t_1][t_2]$ اختصار لـ~$\Atom{\eq}{t_1,t_2}$. ويُختصر أيضًا
$\lnot \Atom{\eq}{t_1,t_2}$ إلى~$\eq/[t_1][t_2]$. وإذا كتبنا الصيغة
$(!B \ast !C)$ المؤلّفة من $!B$ و$!C$ بالرابط الثنائي
المحل~$\ast$، فالغالب أن نسقط قوسيها الخارجيين ونقتصر
على~$!B \ast !C$.

\begin{intro}
وفي بعض كتب المنطق لا بد من ورود ال!!{variable}~$x$ في~$!A$ حتى تُعد
\iftag{prvEx}{$\lexists[x][!A]$ }{}%
\iftag{notprvEx,notprvAll}{}{و }%
\iftag{prvAll}{$\lforall[x][!A]$ }{}%
في عداد
\iftag{notprvEx,notprvAll}{!!a{formula}}{!!{formula}s}.
ولسنا نشترط ذلك؛ فترك هذا الشرط لا يخلّ بالبحث، وفيه تيسير له.
\end{intro}

\begin{tagblock}{defNot,defOr,defAnd,defIf,defIff,defTrue,defFalse,defEx,defAll}
\begin{defn}
والصيغ التي تدخلها المؤثرات المعرّفة تُردّ إلى أصولها على الوجه الآتي:

\begin{tagenumerate}{defTrue,defFalse,defNot,defOr,defAnd,defIf,defIff,defEx,defAll}

\tagitem{defTrue}{$\ltrue$ اختصار لـ
  \iftag{prvFalse}{$\lnot\lfalse$}{$(!A \lor \lnot !A)$، حيث تكون
    !!{formula}~$!A$ ذرية ثابتة ما}.}{}

\tagitem{defFalse}{$\lfalse$ اختصار لـ
  \iftag{prvTrue}{$\lnot\ltrue$}{$(!A \land \lnot !A)$، حيث تكون
    !!{formula}~$!A$ ذرية ثابتة ما}.}{}

\tagitem{defNot}{$\lnot !A$ اختصار لـ~$!A \lif \lfalse$.}{}

\tagitem{defOr}{$!A \lor !B$ اختصار لـ
  \iftag{prvAnd}{$\lnot(\lnot !A \land \lnot !B)$}{$\lnot !A \lif
    !B$}.}{}

\tagitem{defAnd}{$!A \land !B$ اختصار لـ
  \iftag{prvOr}{$\lnot(\lnot !A \lor \lnot !B)$}{$\lnot (!A \lif
    \lnot !B)$}.}{}

\tagitem{defIf}{$!A \lif !B$ اختصار لـ
  \iftag{prvOr}{$(\lnot !A \lor !B)$}{$\lnot (!A \land \lnot !B)$}.}{}

\tagitem{defIff}{$!A \liff !B$ اختصار لـ~$(!A \lif !B) \land (!B
  \lif !A)$.}{}

\tagitem{defAll}{$\lforall[x][!A]$ اختصار لـ~$\lnot\lexists[x][\lnot !A]$.}{}

\tagitem{defEx}{$\lexists[x][!A]$ اختصار لـ~$\lnot\lforall[x][\lnot !A]$.}{}
\end{tagenumerate}
\end{defn}
\end{tagblock}

وقد تختص اللغة بموضع استعمال، فيجري الاصطلاح على وضع ال!!{predicate}s
وال!!{function}s الثنائية المحل بين الحدين اللذين تتعلق بهما؛ مثل
$t_1 < t_2$ و$(t_1 + t_2)$ في لغة الحساب، و$t_1 \in t_2$ في لغة
نظرية المجموعات. وفي لغة الحساب توضع دالة الخلف، على الاصطلاح الجاري،
\emph{بعد} حجتها:~$t'$. وهذه الرسوم جميعًا اختصارات لا غير؛ وأصولها
الصورية، بالترتيب، هي~$\Atom{\Obj A^2_0}{t_1, t_2}$ و
$\Obj f^2_0(t_1, t_2)$ و$\Atom{\Obj A^2_0}{t_1, t_2}$ و$\Obj f^1_0(t)$.

\begin{defn}[الهوية التركيبية]
للهوية التركيبية بين سلاسل الرموز نضع العلامة~$\ident$؛ فقولنا
$!A \ident !B$ يصدق إذا وفقط إذا كان $!A$ و$!B$ سلسلتين من الرموز لهما
الطول نفسه وتحتويان الرمز نفسه في كل موضع.
\end{defn}

وقد يقع على جانبي الرمز~$\ident$ ما تألّف من وصل سلاسل؛ فمثلًا قولنا
$!A \ident (!B \lor !C)$ معناه أن سلسلة الرموز~$!A$ بعينها تحصل من
وصل قوس فاتح بالسلسلة~$!B$، ثم بالرمز~$\lor$، ثم بالسلسلة~$!C$، ثم
بقوس غالق، على هذا الترتيب. ومن ذلك نعلم أن أول رمز في~$!A$
قوس فاتح، وأن في~$!A$ السلسلة~$!B$ جزءًا يبدأ عند الرمز الثاني،
وأن ما يعقب ذلك الجزء هو~$\lor$، وعلى هذا يُقاس الباقي.

ولما كان بناء الحدود وال!!{formula}s استقرائيًا، مبدؤه عناصر أساسية،
كان سبيل البرهان على خواصها مبدأي الاستقراء الآتيين.

\begin{lem}[\emph{مبدأ الاستقراء على الحدود}]
    \ollabel{lem:trmind}
    لتكن $\Lang L$ لغة رتبة أولى.
    إذا كانت خاصية ما~$P$ بحيث
    %
    \begin{enumerate}
        \item تصدق في كل !!{variable}~$v$،
        %
        \item وتصدق في كل !!{constant}~$a$ من~$\Lang L$،
        %
        \item وتصدق في $f(t_1,\dotsc,t_n)$ متى صدقت في
        $t_1$، و\dots، و$t_n$ وكانت $f$ !!{function} ذات $n$ من
            المحال في~$\Lang L$،
    \end{enumerate}
    (على افتراض أن $t_1$، و\dots، و$t_n$ حدود في~$\Lang{L}$)،
    فإن $P$ تصدق في كل حد في~$\Trm[L]$.
\end{lem}

\begin{prob}
    أثبت \olref[fol][syn][frm]{lem:trmind}.
\end{prob}

\begin{lem}[\emph{مبدأ الاستقراء على ال!!{formula}s}]
    \ollabel{thm:frmind}
    لتكن $\Lang L$ لغة رتبة أولى.
    إذا كانت خاصية ما~$P$ تصدق في جميع ال!!{formula}s الذرية
    وكانت بحيث
    %
    \begin{enumerate}
        \tagitem{prvNot}{تصدق في $\lnot !A$ متى صدقت
            في~$!A$؛}{}
        \tagitem{prvAnd}{تصدق في $(!A \land !B)$
            متى صدقت في $!A$ و~$!B$؛}{}
        \tagitem{prvOr}{تصدق في $(!A \lor !B)$
            متى صدقت في $!A$ و~$!B$؛}{}
        \tagitem{prvIf}{تصدق في $(!A \lif !B)$
            متى صدقت في $!A$ و~$!B$؛}{}
        \tagitem{prvIff}{تصدق في $(!A \liff !B)$
            متى صدقت في $!A$ و~$!B$؛}{}
        \tagitem{prvEx}{تصدق في $\lexists[x][!A]$
            متى صدقت في~$!A$؛}{}
        \tagitem{prvAll}{تصدق في $\lforall[x][!A]$
            متى صدقت في~$!A$؛}{}
  \end{enumerate}
  (على افتراض أن $!A$ و$!B$ !!{formula}s في~$\Lang{L}$)،
  فإن $P$ تصدق في جميع الصيغ في~$\Frm[L]$.
\end{lem}

\end{document}
